\subsubsection{乳头夹与夹具：佩戴安全细节}

乳头夹的原理是压迫组织造成局部缺血，取下时血流回流产生集中的"痛快感"——也就是说，它的体验与风险都建立在"可控的暂时缺血"上。安全细节如下：

\begin{itemize}
\item \textbf{时长}：单次佩戴建议不超过 15-20 分钟，取下休息一段时间后再佩戴。用计时器而不是"凭感觉"——缺血组织的感觉会变迟钝，"感觉还行"不等于"还行"；
\item \textbf{观察信号}：夹具远端（乳头顶端）由红转紫、转白，温度明显下降，或出现麻木（痛感消失而非减弱）——任何一项出现，立即取下；
\item \textbf{松紧选择}：从可调式（蝴蝶夹、哑铃夹）的松档开始；固定式金属夹力度大，对新手不友好。佩戴前先揉搓乳头使其充血，受度与舒适度都会更好；
\item \textbf{取下时刻}：血流回流的"风暴"在取下后数秒达到峰值——提前深呼吸、抓稳对方的手，比事后安慰更有用；取下后不要立即冰敷或热敷，让血管反应自然消退；
\item \textbf{禁忌与清洁}：服用抗凝药或有凝血障碍者不适合此类玩法；糖尿病周围神经病变者对"过量"的感知下降，需由同伴协助观察；哺乳期避免强力负压与长时间夹闭。海绵、橡胶与皮革包覆类夹具多孔、难以彻底消毒，宜专人专用。
\end{itemize}
